\section{Results}\label{sec:results}

All tables and figures below are generated from the processed results registry,
as are the values quoted in the running text.

\subsection{Correctness checks}\label{sec:correctness}

All forty-two tests of Appendix~\ref{app:tests} pass, and three checks run over
the results themselves. Every certificate recomputes from its saved Monte-Carlo
estimate, KL and $n_{\mathrm b}$, reproducing all $115$ recorded values to
$10^{-9}$, and this now gates aggregation. Every saved partition recomputes from
its recorded seed and matches index for index, so no reported run used a
label-dependent split whatever its configuration says. And each run asserts
before training that its loaders carry that partition, so the indices checked are
the indices consumed. Each reported $0.99$ confidence is \emph{per certificate,
per run}: a mean over seeds describes several such certificates and is not
itself a $0.99$ bound, which would need a further union over the rows. It is
conditional in one further respect, stated in Section~\ref{sec:confidence}: the
union over prior scales assumes a candidate grid fixed independently of the
bound set, which we argue for but cannot demonstrate. Every $0.99$ below carries
that assumption.

\subsection{Reproduction and main results}\label{sec:main}

\begin{sidewaystable}\centering
\caption{Headline cells (mean~$\pm$~SD over five seeds). Each run's certificate is
a $0$--$1$ PAC-Bayes-kl bound on the Gibbs risk holding at confidence $0.99$; the
column is the descriptive mean of five such per-run certificates and is not
itself a $0.99$ bound, which would need a further union over the five. The two
test errors are descriptive deployment metrics and carry no guarantee.
$n_{\mathrm b}$ is the bound-set size actually used, which halves for a learnt
prior.}\label{tab:main}
\begin{tabular}{llll r r r r r r}
\toprule
Dataset & Model & Prior & Obj. & $n_{\text{bound}}$ & Emp.\ Gibbs & KL$/n_{\text{bound}}$ & \textbf{Certificate} & Gibbs test & Post.mean test \\
\midrule
CIFAR-10 & CNN & learnt & $f_{\text{bbb}}$ & 25000 & $0.0767\pm0.0027$ & $3.35041$ & $0.9801\pm0.0024$ & $0.2412\pm0.0103$ & $0.2162\pm0.0054$ \\
CIFAR-10 & CNN & learnt & $f_{\text{classic}}$ & 25000 & $0.3174\pm0.0086$ & $0.01169$ & $0.3932\pm0.0084$ & $0.2902\pm0.0091$ & $0.2615\pm0.0077$ \\
CIFAR-10 & CNN & learnt & $f_{\text{quad}}$ & 25000 & $0.3003\pm0.0084$ & $0.03373$ & $0.4277\pm0.0074$ & $0.2802\pm0.0086$ & $0.2525\pm0.0077$ \\
CIFAR-10 & CNN & data-free & $f_{\text{quad}}$ & 50000 & $0.9203\pm0.0006$ & $0.00029$ & $0.9297\pm0.0005$ & $0.8978\pm0.0020$ & $0.9000\pm0.0000$ \\
CIFAR-10 & cnn13 & learnt & $f_{\text{quad}}$ & 25000 & $0.2416\pm0.0033$ & $0.04760$ & $0.3881\pm0.0038$ & $0.2202\pm0.0055$ & $0.1791\pm0.0027$ \\
CIFAR-10 & cnn15 & learnt & $f_{\text{quad}}$ & 25000 & $0.3134\pm0.0524$ & $0.05188$ & $0.4722\pm0.0487$ & $0.2847\pm0.0462$ & $0.2265\pm0.0459$ \\
Fashion-MNIST & CNN & learnt & $f_{\text{quad}}$ & 30000 & $0.1069\pm0.0014$ & $0.01411$ & $0.1674\pm0.0048$ & $0.0931\pm0.0024$ & $0.0866\pm0.0011$ \\
Fashion-MNIST & FCN & learnt & $f_{\text{bbb}}$ & 30000 & $0.0913\pm0.0007$ & $0.72021$ & $0.6632\pm0.0045$ & $0.1118\pm0.0016$ & $0.1018\pm0.0017$ \\
Fashion-MNIST & FCN & learnt & $f_{\text{classic}}$ & 30000 & $0.1387\pm0.0013$ & $0.00087$ & $0.1582\pm0.0017$ & $0.1204\pm0.0028$ & $0.1136\pm0.0024$ \\
Fashion-MNIST & FCN & learnt & $f_{\text{quad}}$ & 30000 & $0.1349\pm0.0011$ & $0.00526$ & $0.1746\pm0.0024$ & $0.1186\pm0.0030$ & $0.1109\pm0.0035$ \\
Fashion-MNIST & FCN & data-free & $f_{\text{quad}}$ & 60000 & $0.2461\pm0.0018$ & $0.08728$ & $0.4480\pm0.0011$ & $0.2254\pm0.0062$ & $0.1891\pm0.0056$ \\
MNIST & CNN & learnt & $f_{\text{quad}}$ & 30000 & $0.0201\pm0.0006$ & $0.00414$ & $0.0369\pm0.0013$ & $0.0108\pm0.0011$ & $0.0095\pm0.0009$ \\
MNIST & FCN & learnt & $f_{\text{bbb}}$ & 30000 & $0.0138\pm0.0003$ & $0.24945$ & $0.2658\pm0.0065$ & $0.0187\pm0.0007$ & $0.0149\pm0.0006$ \\
MNIST & FCN & learnt & $f_{\text{classic}}$ & 30000 & $0.0373\pm0.0006$ & $0.00022$ & $0.0455\pm0.0008$ & $0.0240\pm0.0009$ & $0.0192\pm0.0010$ \\
MNIST & FCN & learnt & $f_{\lambda}$ & 30000 & $0.0294\pm0.0005$ & $0.02032$ & $0.0777\pm0.0013$ & $0.0211\pm0.0003$ & $0.0170\pm0.0005$ \\
MNIST & FCN & learnt & $f_{\text{quad}}$ & 30000 & $0.0328\pm0.0006$ & $0.00625$ & $0.0580\pm0.0011$ & $0.0219\pm0.0006$ & $0.0176\pm0.0003$ \\
MNIST & FCN & data-free & $f_{\text{quad}}$ & 60000 & $0.1199\pm0.0020$ & $0.13763$ & $0.3496\pm0.0028$ & $0.0927\pm0.0033$ & $0.0570\pm0.0018$ \\
\bottomrule
\end{tabular}

\end{sidewaystable}

\begin{figure}[h]\centering
\includegraphics[width=0.8\linewidth]{fig_certificate_anatomy.pdf}
\caption{Certificate anatomy: the empirical Gibbs risk and the PAC-Bayes
complexity slack (KL + confidence terms) for the learnt-prior $f_{\text{quad}}$
cells.}\label{fig:anatomy}
\end{figure}

On MNIST with a learnt prior, the tightest $0$--$1$ certificate is obtained by
$f_{\text{classic}}$ ($\CertMnistFcnFclassicLearnt$), followed by $f_{\text{quad}}$
($\CertMnistFcnFquadLearnt$), $f_{\lambda}$ ($\CertMnistFcnFlambLearnt$) and
$f_{\text{bbb}}$ ($\CertMnistFcnBbbLearnt$); all are non-vacuous and reproduce the
qualitative finding of the reference that bound-minimising objectives yield far
tighter certificates than the accuracy-oriented $f_{\text{bbb}}$ baseline, whose
certificate is loose precisely because it carries the largest KL term
(Table~\ref{tab:main}). We do \emph{not} observe the
$f_{\text{quad}}\lesssim f_{\lambda}<f_{\text{classic}}$ ordering sometimes
quoted: the certificate is the same PAC-Bayes-kl inversion for every objective, so
the ordering reflects only which posterior each produces, and here
$f_{\text{classic}}$ yields the lowest-KL one. With the hyper-parameters fixed
this is a comparison of these runs, not a general ordering of the objectives.
Consistently, $f_{\text{bbb}}$ attains the best stochastic test error on MNIST
($\StochMnistFcnBbbLearnt$ against $\StochMnistFcnFquadLearnt$ for $f_{\text{quad}}$)
yet the loosest certificate, because its objective is not calibrated to the
reported bound.

Figure~\ref{fig:anatomy} shows where the looseness comes from. For the tight
learnt-prior cells the certificate sits only a little above the
empirical Gibbs risk: on the MNIST FCN $f_{\text{classic}}$ cell the complexity slack
$\mathrm{KL}/n_{\mathrm b}$ is around $0.0002$, so the bound is almost entirely
empirical risk. As the posterior is allowed a larger KL the slack grows:
$f_{\text{quad}}$ and $f_{\lambda}$ trade slightly lower empirical risk for a
noticeably larger KL, which is why they come out looser than $f_{\text{classic}}$
at similar accuracy, and $f_{\text{bbb}}$ is the extreme, its KL dominating the
certificate even though its empirical risk is the smallest of the four.
Tightness therefore depends on the empirical risk and on the KL cost of reaching
it, not on accuracy alone. That also explains the cross-dataset behaviour: the
empirical-risk floor rises with difficulty (Section~\ref{sec:erroranalysis} shows
where), so even a well-controlled KL cannot pull Fashion-MNIST or CIFAR-10 into
the MNIST range.

\paragraph{Comparison with the published numbers.} Table~\ref{tab:repro} sets our
values against the reference's for the four MNIST cells named as reproduction
targets in Section~\ref{sec:rqs}, using the same-definition columns recorded in
the reference file of Section~\ref{sec:compute}. All four reproduce on the
test-error metric: the stochastic test error differs from the published value by
between $\ReproStochDeltaMin$ and $\ReproStochDeltaMax$, inside the
one-percentage-point tolerance, and
the sign of the difference varies between cells rather than pointing one way, as
it would if our posteriors were systematically better or worse trained. That is
evidence about predictive behaviour, not about each component: the test error is
end-to-end, and the component claims rest on the tests of
Appendix~\ref{app:tests}.

\begin{table}[h]\centering\small
\caption{Reproduction of the four MNIST cells named in Section~\ref{sec:rqs},
against the published values of \textcite{perezortiz2021}. Ours are means over
five seeds. The published numbers were obtained with $m=150{,}000$ Monte-Carlo
samples and a per-cell grid search; ours with $m=2{,}000$ and one fixed
configuration.}\label{tab:repro}
\begin{tabular}{@{}lll rr r rr r@{}}
\toprule
& & & \multicolumn{3}{c}{Certificate} & \multicolumn{3}{c}{Stochastic test error} \\
\cmidrule(lr){4-6}\cmidrule(lr){7-9}
Model & Prior & Obj. & Published & Ours & $\Delta$ & Published & Ours & $\Delta$ \\
\midrule
FCN & data-free & $f_{\text{quad}}$ & $0.3155$ & $0.3496$ & $+0.0341$ & $0.0951$ & $0.0927$ & $-0.0024$ \\
FCN & learnt & $f_{\text{quad}}$ & $0.0279$ & $0.0580$ & $+0.0301$ & $0.0204$ & $0.0219$ & $+0.0015$ \\
FCN & learnt & $f_{\lambda}$ & $0.0354$ & $0.0777$ & $+0.0423$ & $0.0178$ & $0.0211$ & $+0.0033$ \\
CNN & learnt & $f_{\text{quad}}$ & $0.0155$ & $0.0369$ & $+0.0214$ & $0.0127$ & $0.0108$ & $-0.0019$ \\
\bottomrule
\end{tabular}

\end{table}

The certificates do not meet the same tolerance: ours are looser in all four
cells, by between $\ReproCertDeltaMin$ and $\ReproCertDeltaMax$. The Monte-Carlo sweep of
Section~\ref{sec:mc} measures how much of that is budget: re-drawing at the
reference's $m=150{,}000$ takes the MNIST FCN $f_{\text{quad}}$ seed-0 checkpoint
from $0.0572$ to $0.0420$, removing a little over half of that checkpoint's
$0.0293$ gap. The remainder is unexplained. The reproduction therefore succeeds on the
predictive metrics and is partial on the certificate, with one protocol difference
measured and the rest identified but not bounded.

The standard deviations in Table~\ref{tab:main} are small relative to the gaps the
conclusions rest on: the three bound-minimising MNIST learnt-prior certificates
vary by one to two thousandths across the five seeds, whereas the gap between
objectives, or between a learnt and a random prior, is far larger. That the
spread is measurable at all is the seed correction: in the original pipeline the
split and the initialisation were reset to fixed constants regardless of the
command-line seed, so every ``seed'' produced an identical run. With five streams
genuinely driving the split, the initialisation, the loader order, the
Monte-Carlo draws and the evaluation, the spread we report is real, and small
enough that the orderings we draw conclusions from are not artefacts of one lucky
run. It stays separate from the PAC-Bayes confidence, which is about one trained
model rather than variation between runs.

\subsection{Prior comparison}\label{sec:priorcmp}

\begin{figure}[h]\centering
\includegraphics[width=0.7\linewidth]{fig_prior_comparison.pdf}
\caption{Learnt vs random prior certificate, shown only for cells where both
priors were run (including the CIFAR-10 control).}\label{fig:prior}
\end{figure}

The prior changes the KL a useful posterior must pay, and moves certificate
tightness a long way in every paired control (Figure~\ref{fig:prior}). On the
MNIST FCN $f_{\text{quad}}$ cell, going from a random to a learnt prior tightens
the certificate from $\CertMnistFcnFquadRand$ to $\CertMnistFcnFquadLearnt$ while
$\mathrm{KL}(Q\|P)$ falls from $\KlMnistFcnFquadRand$ to
$\KlMnistFcnFquadLearnt$ nats and the empirical Gibbs risk from
$\EmpMnistFcnFquadRand$ to $\EmpMnistFcnFquadLearnt$, so accuracy improves
alongside the certificate. The same direction holds on Fashion-MNIST
($\CertFashionmnistFcnFquadRand\!\to\!\CertFashionmnistFcnFquadLearnt$).

Which of the two terms does the work is a separate question, and the paired cells
answer it differently. The quantity to compare is not the raw KL but
$\mathrm{KL}/n_{\mathrm b}$, since $n_{\mathrm b}$ halves when the prior is
learnt. On MNIST the complexity term falls by $0.131$ against $0.087$ for the
empirical risk, so the complexity route is the larger. On Fashion-MNIST the two
are comparable and the ranking reverses: $0.082$ against $0.111$. On CIFAR-10 the
complexity term \emph{rises} by $0.033$, because the random-prior posterior barely
leaves its initialisation and so pays almost no KL, and the whole of the gain is
the empirical risk falling from $\EmpCifaronezeroCnnFquadRand$ to
$\EmpCifaronezeroCnnFquadLearnt$. The learnt prior tightens the certificate in
all three, then, but not by one mechanism: across these three paired cells the
dominant contribution shifts from the complexity term to the empirical risk. We
do not attribute that shift to difficulty, since architecture and budget change
along the same sequence.
The comparison is between two whole protocols rather than a pure prior-family
intervention: the random prior is independent of all of $S$ and so uses
$n_{\mathrm b}=|S|$, while the learnt prior uses the smaller
$n_{\mathrm b}=|S\setminus S_0|$ --- and still attains the tighter certificate
\emph{despite} that handicap. CIFAR-10 is the clearest case: a random prior
leaves the $f_{\text{quad}}$ certificate at $\CertCifaronezeroCnnFquadRand$,
whereas the learnt prior moves it to $\CertCifaronezeroCnnFquadLearnt$, formally
non-vacuous but still loose.

\paragraph{Architecture.} Where a CNN and an FCN are compared on the same dataset,
the convolutional model gives the tighter certificate: on MNIST the $f_{\text{quad}}$
learnt-prior certificate is $\CertMnistCnnFquadLearnt$ (CNN) versus
$\CertMnistFcnFquadLearnt$ (FCN), and on Fashion-MNIST
$\CertFashionmnistCnnFquadLearnt$ versus $\CertFashionmnistFcnFquadLearnt$.

\paragraph{Depth on CIFAR-10.} The reference obtains its tighter CIFAR-10
certificates with deeper networks than the nine-layer one used above, so we ran
its 13- and 15-layer architectures at the same five seeds and the same fixed
configuration. Within this $f_{\text{quad}}$ depth sweep the mean certificate
improves from nine to thirteen layers and worsens at fifteen:
$\CertCifaronezeroCnnFquadLearnt$, $\CertCifaronezeroCnnonethreeFquadLearnt$ and
$\CertCifaronezeroCnnonefiveFquadLearnt$ respectively. The 13-layer network also
improves the deployed error ($\PostmeanCifaronezeroCnnonethreeFquadLearnt$ versus
$\PostmeanCifaronezeroCnnFquadLearnt$), and the gain concentrates where the error
is: the per-class Gibbs risk of Section~\ref{sec:erroranalysis} falls from $0.448$
to $0.340$ on \emph{cat}, the hardest class. It is also the tightest CIFAR-10
certificate in the study, narrowly ahead of $f_{\text{classic}}$ at nine layers
($\CertCifaronezeroCnnFclassicLearnt$), and so the best of the three depths tested
under $f_{\text{quad}}$ and this configuration.

The 15-layer network does not continue the trend, for an instructive reason. Two
of its five seeds barely left the chance-level plateau, their prior networks
finishing at $0.866$ and $0.899$ error on the $S_0$ they were fitted to, against
a chance level of $0.900$: the second is chance to three decimals and the first
sits inside the criterion's five-point tolerance, both far from the at most
$0.228$ every other learnt prior reaches. The posterior is then initialised at an
essentially untrained prior, so the cell is not an instance of the protocol. Both are
registered as \texttt{not\_converged} and excluded by the criterion of
Section~\ref{sec:compute}; the remaining three give
$\CertCifaronezeroCnnonefiveFquadLearnt$, looser than the 13-layer result and far
more variable. That mean is conditional on convergence, so the failure rate is
part of the result: at fifteen layers this protocol trains $3$ of $5$ seeds,
against $5$ of $5$ everywhere else in the study. Twelve convolutional layers with
no normalisation, no warm-up and one fixed learning rate is past what it trains
reliably --- a limit of the fixed reduced-compute configuration rather than of the
method, and the price of holding one configuration across every cell.

\paragraph{Across datasets.} The learnt-prior $f_{\text{quad}}$ certificate rises
from MNIST ($\CertMnistFcnFquadLearnt$, FCN) through Fashion-MNIST
($\CertFashionmnistFcnFquadLearnt$, FCN) to CIFAR-10
($\CertCifaronezeroCnnFquadLearnt$, 9-layer CNN). The pattern is descriptive,
since architecture and compute budget change along it;
Section~\ref{sec:difficulty} is the controlled test of difficulty.

\subsection{Monte-Carlo sensitivity}\label{sec:mc}

\begin{figure}[h]\centering
\includegraphics[width=0.7\linewidth]{fig_mc_convergence.pdf}
\caption{Final certificate versus the Monte-Carlo budget $m$ on a fixed checkpoint
with the Monte-Carlo estimate genuinely re-drawn at each budget, so both the
estimate and the finite-MC slack are free to move.}\label{fig:mc}
\end{figure}

The comparison above attributed our looser certificates partly to the Monte-Carlo
budget; this sweep measures that. Holding the posterior fixed at a saved
checkpoint and re-drawing at each budget, the certificate tightens monotonically
as $m$ grows and flattens well before the reference's budget
(Figure~\ref{fig:mc}). On the MNIST FCN $f_{\text{quad}}$ checkpoint it falls from
$0.0654$ at $m=1{,}000$ to $0.0572$ at our $m=2{,}000$ and $0.0420$ at the
reference's $m=150{,}000$ --- all the \emph{same} seed-0 checkpoint, which is why
the $m=2{,}000$ entry is not the five-seed mean $\CertMnistFcnFquadLearnt$ of
Table~\ref{tab:main}. Against the published $0.0279$, the re-draw recovers
$0.0152$ of this checkpoint's $0.0293$ gap. The rest is unexplained; the unmatched
grid search and epoch budget are plausible contributors we did not test
separately.

Two checks confirm that what moves is the slack and not the estimate. The raw
Monte-Carlo estimate is stable across the whole range ($0.01983$ at $m=1{,}000$
against $0.01984$ at $m=150{,}000$), as an unbiased estimator should be; and the
same curve computed analytically, holding that estimate fixed, agrees with the
measured one to within $2.2\times10^{-5}$ at every budget. Both curves are in the
artefact. A re-draw could have
moved the certificate either way; here it fell monotonically at every step, so
for this checkpoint and these draws the smaller $m$ was the conservative one.

\subsection{Small-data behaviour}\label{sec:smalldata}

\begin{table}[h]\centering\small
\caption{MNIST FCN $f_{\text{quad}}$/learnt versus training-data fraction, with the
corrected bound-set size $n_{\mathrm b}$ (mean~$\pm$~SD over seeds).}\label{tab:smalldata}
\begin{tabular}{r r r r r}
\toprule
\% data & $n_{\text{posterior}}$ & $n_{\text{bound}}$ & Emp.\ Gibbs & \textbf{Certificate} \\
\midrule
100\% & 60000 & 30000 & $0.0328\pm0.0006$ & $0.0580\pm0.0011$ \\
50\% & 30000 & 15000 & $0.0420\pm0.0024$ & $0.0777\pm0.0049$ \\
20\% & 12000 & 6000 & $0.0624\pm0.0034$ & $0.1196\pm0.0040$ \\
10\% & 6000 & 3000 & $0.0771\pm0.0035$ & $0.1520\pm0.0065$ \\
\bottomrule
\end{tabular}

\end{table}

\begin{figure}[h]\centering
\includegraphics[width=0.7\linewidth]{fig_smalldata.pdf}
\caption{Certificate versus training-data fraction with the corrected
$n_{\mathrm b}$ and across-seed error bars.}\label{fig:smalldata}
\end{figure}

The certificate loosens monotonically as the training set shrinks
(Table~\ref{tab:smalldata} and Figure~\ref{fig:smalldata}): from
$\CertMnistFcnFquadLearnt$ at full data to
$\CertMnistFcnFquadLearntPtFiveZero$ ($50\%$),
$\CertMnistFcnFquadLearntPtTwoZero$ ($20\%$) and
$\CertMnistFcnFquadLearntPtOneZero$ at $10\%$ of the data
($n_{\mathrm b}=3{,}000$). The certificate increases monotonically as the data
fraction falls, and each step exceeds the across-seed spread at that setting. The small-data runs show most directly why the
sample-count audit matters. The certificate denominator is the bound-set size
$n_{\mathrm b}$, which at $10\%$ data is $3{,}000$, whereas taking it to be the
size of the full training set would instead give $60{,}000$. Using the full
dataset size understates the complexity term and reports a spuriously tight
certificate at small $n$: on these same five checkpoints the wrong denominator
returns $0.092$ where the correct one returns
$\CertMnistFcnFquadLearntPtOneZero$, hiding most of the degradation. The bound
stays non-vacuous down to $10\%$ of the data, so the degradation is graceful, but
it is far from the near-flat behaviour a mis-counted $n$ would suggest.

\subsection{Controlled difficulty}\label{sec:difficulty}

\begin{figure}[h]\centering
\includegraphics[width=0.7\linewidth]{fig_difficulty.pdf}
\caption{Empirical risk, $\mathrm{KL}/n_{\mathrm b}$ and certificate versus
label-noise level for the MNIST FCN $f_{\text{quad}}$/learnt cell (architecture,
sample size and compute held fixed).}\label{fig:difficulty}
\end{figure}

The cross-dataset trend above could not separate difficulty from the architecture
and budget that change alongside it; this study removes that confound. With
architecture, sample size and compute held fixed and label noise injected into the
MNIST FCN $f_{\text{quad}}$ learnt-prior cell, the certificate rises monotonically
with the noise level (Figure~\ref{fig:difficulty}), from
$\CertMnistFcnFquadLearnt$ at no noise to
$\CertMnistFcnFquadLearntLnOneZero$, $\CertMnistFcnFquadLearntLnTwoZero$ and
$\CertMnistFcnFquadLearntLnFourZero$ at $10$, $20$ and $40\%$ noise. Successive
noise levels are separated by far more than the across-seed spread
(Figure~\ref{fig:difficulty}). The increase is
driven mainly by the empirical-risk term, since corrupted labels raise the error
floor that any posterior must pay on $S_{\mathrm b}$, with a smaller rise in
$\mathrm{KL}/n_{\mathrm b}$. Because
every other factor is held fixed, this is the one place we attribute a difficulty
effect to the intervention itself. The certified quantity here is the Gibbs risk
on the \emph{noisy-label} distribution the model is trained and evaluated
against, not on clean labels.

\subsection{The certified predictor versus the deployed one}\label{sec:deploygap}

Table~\ref{tab:main} also lets us compare the predictor the certificate covers
against the one a user would deploy. The posterior-mean test error is below the
stochastic (Gibbs) error in every learnt-prior cell, though not universally: on
the CIFAR-10 data-free-prior cell, where neither predictor has learnt much, the
ordering reverses. The finite ensemble is usually below the stochastic error but
is not uniformly below the posterior mean. On the MNIST FCN $f_{\text{quad}}$
cell the stochastic error is $\StochMnistFcnFquadLearnt$ against a posterior mean
of $\PostmeanMnistFcnFquadLearnt$, and the same direction holds on Fashion-MNIST
and CIFAR-10 with a learnt prior. The gap is consistent across the learnt-prior
cells and not negligible, at roughly a fifth of the stochastic error on the MNIST
cell, and by Section~\ref{sec:scope} it does not transfer the certificate to the
posterior mean, so a user needing a guarantee for the deterministic network does
not have one here. Closing that gap is the clearest theoretical extension of the
project.

\subsection{Where the error falls}\label{sec:erroranalysis}

The certificate bounds an aggregate risk, and it is informative to see how the
error distributes across classes, since that is what the empirical-risk term
integrates over. Figure~\ref{fig:confusion} shows the posterior-mean confusion
matrices on the two harder datasets. On Fashion-MNIST the error concentrates in
the upper-body garments, the worst class being \emph{shirt}, confused mainly with
\emph{pullover}, \emph{coat} and \emph{T-shirt}, while the footwear classes are
almost perfectly separated; on CIFAR-10 it sits in \emph{cat}, \emph{dog} and
\emph{bird}, with the rigid vehicle classes much easier. Those matrices describe
the posterior-mean predictor on the test
set, which is not the object the certificate bounds, so we also decomposed the
certified quantity itself: the Monte-Carlo Gibbs risk on $S_{\mathrm b}$,
accumulated by true class over the same $m=2{,}000$ draws that produce the
certificate, on the seed-$0$ checkpoint of each cell. What is decomposed is the
Monte-Carlo average $\bar{R}$, before the finite-sample correction
of~\eqref{eq:mc}, so the per-class values below sit on that scale rather than on
the scale of the corrected $\overline{R}^{+}$ quoted elsewhere. Its class average reproduces the
run's own aggregate estimate to within $10^{-4}$, which is the check that the
decomposition is of the right quantity.

The certified risk concentrates in the same places, and more sharply: on
Fashion-MNIST the Gibbs risk on \emph{shirt} is $0.293$ against a dataset-wide
$0.108$ on the same scale, and on CIFAR-10 it is \emph{cat} $0.448$,
\emph{dog} $0.383$ and \emph{bird} $0.366$ against $0.264$ for the nine-layer
network. The worst class carries close to three times the average risk the bound
integrates over on Fashion-MNIST and about $1.7$ times it on CIFAR-10. That is
what the empirical-risk term of Section~\ref{sec:main} is made of, and why a
well-controlled KL cannot pull the harder datasets into the MNIST range. It also
marks a boundary of the method: the certificate is one number over all classes, so
a deployment that cares about the worst class is not served by it.

\begin{figure}[h]\centering
\begin{subfigure}{0.49\linewidth}\centering
\includegraphics[width=\linewidth]{fig_confusion_fashion-mnist.pdf}
\caption{Fashion-MNIST}\end{subfigure}
\begin{subfigure}{0.49\linewidth}\centering
\includegraphics[width=\linewidth]{fig_confusion_cifar10.pdf}
\caption{CIFAR-10}\end{subfigure}
\caption{Posterior-mean confusion matrices (row-normalised) on the two harder
datasets. The error concentrates in visually similar classes, which raises the
empirical Gibbs risk that the certificate is built on.}\label{fig:confusion}
\end{figure}
